% !TeX root = tcolorbox.tex
% include file of tcolorbox.tex (manual of the LaTeX package tcolorbox)
\clearpage
\section{记录}\label{sec:recording}%
\tcbset{external/prefix=external/recording_}%
本宏包提供了若干宏和选项，用于在编译期间进行\emph{记录}。
这是通过 \LaTeX\ 文件操作实现的：把某些数据保存到文件中，供后续使用。
主要应用场景见 \zvref{sec:recording-exercises}，其中记录了示例解答的信息，
并在 \zvref{sec:recording-solutions} 中再次读入。

\subsection{宏}\label{sec:recording-makros}
\begin{docCommand}[doc new=2014-11-28]{tcbstartrecording}{\oarg{file name}}
  打开由 \meta{文件名}指定的文件，用于写入记录。
  默认文件名为 |\jobname.records|。
  应用示例见 \zvref{sec:recording-exercises}。
  \begin{marker}
  在某些情况下，缺失可选参数可能导致解析问题。
  若发生这种情况（或仅作预防），请使用
\begin{dispListing}
\tcbstartrecording\relax
\end{dispListing}
  或
\begin{dispListing}
\tcbstartrecording[\jobname.records]
\end{dispListing}
  \end{marker}
\end{docCommand}

\begin{docCommand}[doc new=2014-11-28]{tcbrecord}{\marg{content}}
  把任意 \meta{内容}记录到记录文件中。
  \refCom{tcbrecord} 以 |\immediate\write| 实现。
  须先调用 \refCom{tcbstartrecording}；否则 \refCom{tcbrecord} 会被静默忽略。
\begin{dispListing}
\tcbrecord{\string\solution{\thetcbcounter}{solutions/exercise-\thetcbcounter.tex}}
\end{dispListing}
\end{docCommand}

\begin{docCommand}[doc new=2014-11-28]{tcbstoprecording}{}
  关闭此前由 \refCom{tcbstartrecording} 打开的当前记录文件。
\end{docCommand}

\begin{docCommand}[doc new=2014-11-28]{tcbinputrecords}{\oarg{file name}}
  打开由 \meta{文件名}指定的文件，通过 |\input| 读取其中的记录。
  默认文件名是上一次用于保存的记录文件名。
  须先调用 \refCom{tcbstoprecording}。
\end{docCommand}


\subsection{选项}\label{sec:recording-options}
\begin{docTcbKey}[][doc new=2014-11-28]{record}{=\meta{content}}{样式，无默认值}
  把任意 \meta{内容}记录到记录文件中，见 \refCom{tcbrecord}。
  本键可多次使用，以写入多行。
  \begin{dispListing}
  record={\string\solution{\thetcbcounter}{solutions/exercise-\thetcbcounter.tex}}
  \end{dispListing}
\end{docTcbKey}

\begin{docTcbKey}[][doc new=2014-11-28]{no recording}{}{}
  在当前分组内禁用 \refCom{tcbrecord} 与 \refKey{/tcb/record}。
\end{docTcbKey}


\clearpage
\subsection{示例：练习题}\label{sec:recording-exercises}
下面的应用示例创建练习题及其对应解答。
每一对都生成在单个 |tcolorbox| 内，其中解答位于 \refCom{tcblower} 之后。
对于每个示例，解答部分通过 \refKey{/tcb/savelowerto} 保存到文件，
并通过 \refKey{/tcb/record} 记录此次保存。为展示其可能性，
第二道练习题没有解答。最后，解答在 \zvref{sec:recording-solutions} 中被读入。

\inputpreamblelisting{L}

\begin{dispListing*}{breakable,before upper=}
\tcbstartrecording\relax

\begin{exercise}
  Compute the derivative of the following function:
  \begin{equation*}
    f(x)=\sin((\sin x)^2)
  \end{equation*}
\tcblower
  The derivative is:
  \begin{align*}
    f'(x) &= \left( \sin((\sin x)^2) \right)'
    =\cos((\sin x)^2) 2\sin x \cos x.
  \end{align*}
\end{exercise}

\begin{exercise}[no solution]
  It holds:
  \begin{equation*}
    \frac{d}{dx}\left(\ln|x|\right) = \frac{1}{x}.
  \end{equation*}
\end{exercise}

\begin{exercise}
  Compute the derivative of the following function:
  \begin{equation*}
    f(x)=(\sin(\sin x))^2
  \end{equation*}
\tcblower
  The derivative is:
  \begin{align*}
    f'(x) &= \left( (\sin(\sin x))^2 \right)'
    =2\sin(\sin x)\cos(\sin x)\cos x.
  \end{align*}
\end{exercise}

\begin{exercise}
  Compute the derivative of the following function:
  \begin{equation*}
    f(x)=\sqrt{x^3-6x^2+2x}
  \end{equation*}
\tcblower
  The derivative is:
  \begin{align*}
    f'(x) &= \left( \sqrt{x^3-6x^2+2x} \right)'
    = \frac{3x^2-12x+2}{2\sqrt{x^3-6x^2+2x}}.
  \end{align*}
\end{exercise}

\begin{exercise}
  Compute the derivative of the following function:
  \begin{equation*}
    f(x)=\left(\frac{2+3x}{1-2x}\right)^3
  \end{equation*}
\tcblower
  The derivative is:
  \begin{align*}
    f'(x) &= \left( \left(\frac{2+3x}{1-2x}\right)^3 \right)'
    = 3 \left(\frac{2+3x}{1-2x}\right)^2 \frac{(1-2x)3-(2+3x)(-2)}{(1-2x)^2}
    = \frac{21(2+3x)^2}{(1-2x)^4}.
  \end{align*}
\end{exercise}

\begin{exercise}
  Compute the derivative of the following function:
  \begin{equation*}
    f(x)=\frac{\cos x}{(\tan 2x)^2}
  \end{equation*}
\tcblower
  The derivative is:
  \begin{align*}
    f'(x) &= \left( \frac{\cos x}{(\tan 2x)^2} \right)'
    = \left( \frac{\cos x (\cos 2x)^2}{(\sin 2x)^2} \right)'\\
    &= \frac{(\sin 2x)^2 [(-\sin x)(\cos 2x)^2+(\cos x)4\cos 2x (-\sin 2x)]
       - \cos x (\cos 2x)^2 4\sin 2x \cos 2x}{(\sin 2x)^4}\\
    &= -\frac{\cos(2x) [\sin x \sin 2x \cos 2x+ 4\cos x(\sin 2x)^2
       + 4 \cos x (\cos 2x)^2]}{(\sin 2x)^3}\\
    &= -\frac{\cos(2x) [\sin x \sin 2x \cos 2x+ 4\cos x]}{(\sin 2x)^3}.
  \end{align*}
\end{exercise}

\begin{exercise}
  Compute the derivative of the following function:
  \begin{equation*}
    f(x)=\cos((2x^2+3)^3)
  \end{equation*}
\tcblower
  The derivative is:
  \begin{align*}
    f'(x) &= \left( \cos((2x^2+3)^3) \right)'
    =-\sin((2x^2+3)^3) 3(2x^2+3)^2 2\cdot 2x\\
    &=-12x(2x^2+3)^2\sin((2x^2+3)^3).
  \end{align*}
\end{exercise}

\begin{exercise}
  Compute the derivative of the following function:
  \begin{equation*}
    f(x)=(x^2+1)\sqrt{x^4+1}
  \end{equation*}
\tcblower
  The derivative is:
  \begin{align*}
    f'(x) &= \left( (x^2+1)\sqrt{x^4+1} \right)'
    = 2x\sqrt{x^4+1} + \frac{2x^3(x^2+1)}{\sqrt{x^4+1}}.
  \end{align*}
\end{exercise}

\tcbstoprecording
\end{dispListing*}
\tcbusetemp

\subsection{示例：解答}\label{sec:recording-solutions}
此处承接 \zvref{sec:recording-exercises} 给出的示例。
现在读入已保存并记录的解答。

\begin{dispListing*}{breakable,before upper=}
\tcbinputrecords
\end{dispListing*}
\tcbusetemp
